\chapter{État de l'art}
\label{chap:etat-art}

Le chapitre précédent a établi ce que le système doit faire. Celui-ci examine ce que
la communauté scientifique sait déjà faire, et ce qu'elle ne sait pas encore. Nous
parcourons successivement les trois corpus de travaux sur lesquels s'appuie notre
proposition : la classification automatique des stades du sommeil, l'apprentissage
ensembliste, et l'intelligence artificielle explicable. Une quatrième section
s'intéresse aux méthodes de dépistage du syndrome d'apnées obstructives, qui
constituent le second volet du projet. Nous concluons en situant précisément notre
contribution dans ce paysage.

% =============================================================================
\section{La classification automatique des stades du sommeil}
% =============================================================================

\subsection{Les approches fondées sur des caractéristiques expertes}

Jusqu'au milieu des années 2010, la démarche dominante consistait à reproduire le
raisonnement du scoreur humain en deux temps. Un ingénieur du signal concevait
d'abord un jeu de descripteurs censés capturer les éléments décrits par le référentiel
AASM : puissance spectrale dans les bandes delta, thêta, alpha et bêta, entropie de
permutation, paramètres de Hjorth, coefficients d'ondelettes. Un classifieur classique
— machine à vecteurs de support, forêt aléatoire, $k$ plus proches voisins — apprenait
ensuite à associer ces descripteurs aux stades.

Fraiwan \emph{et al.}~\cite{fraiwan2012} illustrent bien cette famille : une analyse
temps-fréquence d'un unique canal EEG alimente une forêt aléatoire. Hassan et
Bhuiyan~\cite{hassan2016} y ajoutent une décomposition modale empirique et une
agrégation bootstrap.

Ces méthodes présentent un avantage considérable, souvent oublié dans les comparaisons
de performance : elles sont \emph{intrinsèquement interprétables}. Lorsqu'une forêt
aléatoire classe une époque en N3 parce que la puissance dans la bande delta dépasse
un seuil, la justification est immédiatement lisible par un clinicien, car elle emploie
le vocabulaire même du référentiel de scorage.

Leur limite est ailleurs. La qualité du résultat dépend entièrement de la pertinence
des descripteurs choisis \emph{a priori}. Un motif discriminant que l'ingénieur n'a pas
anticipé reste invisible au classifieur, quelle que soit sa puissance. C'est
précisément ce plafond que l'apprentissage de représentations vient lever.

\subsection{L'apprentissage de représentations : les réseaux convolutifs}

Un réseau de neurones convolutif~\cite{lecun1998} apprend ses propres filtres au lieu
de les recevoir. Appliqué à un signal unidimensionnel, il extrait des motifs locaux
invariants par translation — ce qui correspond exactement à la nature des éléments
recherchés dans une époque : un fuseau du sommeil ou un complexe K sont des formes
brèves dont la position dans la fenêtre de 30 secondes est indifférente.

Tsinalis \emph{et al.}~\cite{tsinalis2016} ont été parmi les premiers à appliquer un
CNN à l'EEG brut, sans extraction préalable. Chambon \emph{et al.}~\cite{chambon2018}
ont étendu l'approche au cas multivarié, en traitant conjointement EEG, EOG et EMG,
configuration proche de celle que nous adoptons.

Une propriété architecturale mérite d'être soulignée car elle guidera notre conception :
la taille du noyau de convolution détermine l'échelle temporelle observée. Un noyau
court capte les oscillations rapides ; un noyau long, les ondes lentes. Comme les
stades du sommeil se distinguent précisément par des phénomènes d'échelles différentes,
les architectures les plus efficaces emploient plusieurs branches parallèles à
résolutions distinctes. C'est le principe retenu par
DeepSleepNet~\cite{supratak2017} et que nous reprendrons.

\subsection{La modélisation de la dépendance temporelle}

Un scoreur humain ne juge jamais une époque isolément. Le référentiel AASM impose
d'ailleurs des règles de continuité : certaines transitions sont interdites, et le
contexte des époques voisines conditionne l'attribution. Un modèle qui classerait
chaque fenêtre indépendamment se priverait de cette information.

Les réseaux récurrents à mémoire longue~\cite{hochreiter1997} maintiennent un état
interne le long de la séquence et apprennent quelles informations conserver ou
oublier. Leur variante bidirectionnelle parcourt la séquence dans les deux sens, ce
qui autorise une décision informée par le passé comme par l'avenir — situation
réaliste ici, puisque l'enregistrement est analysé après coup et non en temps réel.

DeepSleepNet~\cite{supratak2017} combine ainsi une extraction convolutive
multi-échelle et un LSTM bidirectionnel, architecture devenue la référence du domaine.
TinySleepNet~\cite{supratak2020} en propose une version considérablement allégée pour
les dispositifs embarqués. SleepEEGNet~\cite{mousavi2019} traite le problème comme une
tâche de traduction séquence à séquence.

\subsection{L'attention et les architectures Transformer}

Le mécanisme d'attention~\cite{bahdanau2015} permet à un modèle de pondérer
dynamiquement les éléments d'une séquence selon leur pertinence pour la décision
courante. Les architectures Transformer~\cite{vaswani2017} en font leur unique
primitive de calcul, abandonnant toute récurrence : chaque position peut consulter
directement toutes les autres, ce qui supprime le goulet d'étranglement de l'état
récurrent et autorise une parallélisation complète de l'entraînement.

AttnSleep~\cite{eldele2021} associe une extraction multi-résolution à un module
d'auto-attention. SleepTransformer~\cite{phan2022} construit une architecture
hiérarchique, à l'échelle intra-époque puis inter-époques, et présente un intérêt
particulier pour notre problématique : les auteurs y adjoignent une quantification de
l'incertitude et exploitent les poids d'attention comme support d'interprétation.
XSleepNet~\cite{phan2021} exploite conjointement la représentation temporelle brute et
la représentation temps-fréquence. U-Sleep~\cite{perslev2021}, enfin, adopte une
architecture entièrement convolutive et se distingue moins par sa précision brute que
par sa robustesse : les auteurs valident le modèle sur une quinzaine de cohortes aux
montages hétérogènes.

\subsection{Synthèse comparative}

Le tableau~\ref{tab:etat-art-staging} récapitule ces travaux.

% ─────────────────────────────────────────────────────────────────────────────
% À VÉRIFIER : les exactitudes rapportées ci-dessous doivent être recoupées avec
% les articles originaux avant le dépôt. Les protocoles d'évaluation diffèrent
% (jeux de données, nombre de canaux, validation croisée), ce qui interdit toute
% comparaison directe — le paragraphe qui suit le tableau le précise.
% ─────────────────────────────────────────────────────────────────────────────
\begin{table}[H]
\centering
\caption{Principaux travaux de classification automatique des stades du sommeil}
\label{tab:etat-art-staging}
\small
\begin{tabular}{@{}L{3.3cm}cL{4.4cm}L{2.2cm}c@{}}
\toprule
\textbf{Travail} & \textbf{Année} & \textbf{Approche} & \textbf{Évaluation} &
\textbf{Exactitude} \\
\midrule
Fraiwan \emph{et al.}~\cite{fraiwan2012} & 2012 & Temps-fréquence + forêt aléatoire &
jeu propriétaire & $\approx$ \SI{83}{\percent} \\
\addlinespace[2pt]
Tsinalis \emph{et al.}~\cite{tsinalis2016} & 2016 & CNN sur EEG brut & Sleep-EDF &
$\approx$ \SI{82}{\percent} \\
\addlinespace[2pt]
DeepSleepNet~\cite{supratak2017} & 2017 & CNN multi-échelle + BiLSTM & Sleep-EDF &
$\approx$ \SI{82}{\percent} \\
\addlinespace[2pt]
Chambon \emph{et al.}~\cite{chambon2018} & 2018 & CNN multivarié (EEG, EOG, EMG) &
MASS & $\approx$ \SI{83}{\percent} \\
\addlinespace[2pt]
SleepEEGNet~\cite{mousavi2019} & 2019 & CNN + séquence à séquence & Sleep-EDF &
$\approx$ \SI{84}{\percent} \\
\addlinespace[2pt]
TinySleepNet~\cite{supratak2020} & 2020 & CNN + LSTM allégés & Sleep-EDF &
$\approx$ \SI{85}{\percent} \\
\addlinespace[2pt]
AttnSleep~\cite{eldele2021} & 2021 & Multi-résolution + auto-attention & Sleep-EDF &
$\approx$ \SI{84}{\percent} \\
\addlinespace[2pt]
U-Sleep~\cite{perslev2021} & 2021 & Entièrement convolutif & 15 cohortes &
F1 macro $\approx$ 0,79 \\
\addlinespace[2pt]
SleepTransformer~\cite{phan2022} & 2022 & Transformer hiérarchique & SHHS &
$\approx$ \SI{88}{\percent} \\
\addlinespace[2pt]
XSleepNet~\cite{phan2021} & 2022 & Fusion brut + temps-fréquence & SHHS &
$\approx$ \SI{89}{\percent} \\
\bottomrule
\end{tabular}
\end{table}

Ces valeurs appellent une mise en garde. Elles proviennent de protocoles
d'évaluation différents — jeux de données, nombre de canaux, nombre de classes,
schéma de validation — et \textbf{ne sont pas directement comparables entre elles}.
Un modèle évalué sur Sleep-EDF, cohorte de sujets majoritairement sains, n'affronte
pas la même difficulté qu'un modèle évalué sur SHHS, cohorte âgée et
pathologique~\cite{quan1997}. Le tableau doit donc se lire comme une cartographie des
approches, non comme un classement.

Deux constats se dégagent malgré tout. D'une part, les performances plafonnent depuis
plusieurs années autour de l'accord inter-expert évoqué au
chapitre~\ref{chap:cadre} : le gain marginal d'une architecture nouvelle se compte
désormais en fractions de point. D'autre part, aucune famille architecturale ne domine
uniformément — observation qui motive directement le recours à l'apprentissage
ensembliste.

% =============================================================================
\section{L'apprentissage ensembliste}
% =============================================================================

\subsection{Trois familles de combinaison}

L'idée fondatrice de l'apprentissage ensembliste est qu'un collectif de modèles
imparfaits peut surpasser le meilleur d'entre eux, à condition que leurs erreurs ne
soient pas identiques~\cite{dietterich2000}. Trois familles se distinguent par la
manière dont elles produisent cette diversité.

Le \textbf{bagging}~\cite{breiman1996} entraîne des modèles de même nature sur des
rééchantillonnages aléatoires du jeu d'apprentissage, puis moyenne leurs prédictions.
La forêt aléatoire~\cite{breiman2001} y ajoute un tirage aléatoire des variables à
chaque nœud. La diversité provient donc des données, et l'effet principal est une
réduction de la variance.

Le \textbf{boosting}~\cite{freund1997} procède séquentiellement : chaque modèle est
entraîné à corriger les erreurs de ses prédécesseurs. Sa formulation par descente de
gradient~\cite{friedman2001} a donné les implémentations qui dominent aujourd'hui les
données tabulaires, XGBoost~\cite{chen2016} et LightGBM~\cite{ke2017}. L'effet
principal est ici une réduction du biais.

L'\textbf{empilement}, ou \emph{stacking}~\cite{wolpert1992}, diffère des deux
précédents sur un point essentiel : la diversité ne vient ni des données ni de la
séquence d'entraînement, mais de l'\emph{hétérogénéité des modèles eux-mêmes}. Des
architectures de natures différentes sont entraînées indépendamment, et un second
modèle — le méta-apprenant — apprend à combiner leurs sorties.

\subsection{Le principe de l'empilement}

L'empilement s'organise en deux niveaux, représentés à la
figure~\ref{fig:stacking-principe}. Au premier niveau, $M$ modèles de base produisent
chacun un vecteur de probabilités sur les $C$ classes. Ces vecteurs sont concaténés
pour former un vecteur de méta-caractéristiques de dimension $M \times C$, qui
constitue l'entrée du second niveau.

\begin{figure}[H]
\centering
\begin{tikzpicture}[
    x=1cm,y=1cm,
    bloc/.style   = {rectangle,rounded corners=2pt,draw=hypnoblue!60,fill=hypnolight,
                     align=center,font=\small,minimum height=1.0cm,text width=2.5cm},
    meta/.style   = {rectangle,rounded corners=2pt,draw=hypnored!75,fill=hypnored!8,
                     align=center,font=\small,minimum height=1.0cm,text width=2.8cm},
    ent/.style    = {rectangle,rounded corners=2pt,draw=hypnogray!60,fill=white,
                     align=center,font=\small,minimum height=1.0cm,text width=2.0cm},
    fl/.style     = {-{Latex[length=2.2mm]},hypnoblue!75,line width=0.8pt},
    lab/.style    = {font=\scriptsize\itshape,text=hypnogray},
  ]

  \node[ent]  (x)   at (0,0)      {Époque\\$(5 \times 3000)$};

  \node[bloc] (cnn) at (3.6,2.0)  {CNN};
  \node[bloc] (lstm)at (3.6,0)    {BiLSTM};
  \node[bloc] (tr)  at (3.6,-2.0) {Transformer};

  \node[ent]  (cat) at (7.3,0)    {Concaténation\\$(M \times C)$};
  \node[meta] (ml)  at (10.9,0)   {Méta-apprenant\\\scriptsize régression logistique};
  \node[ent]  (out) at (14.1,0)   {Stade\\prédit};

  \draw[fl] (x) -- (cnn);
  \draw[fl] (x) -- (lstm);
  \draw[fl] (x) -- (tr);

  \draw[fl] (cnn) -- node[lab,above,pos=0.5,yshift=1pt] {$p_1$} (cat);
  \draw[fl] (lstm)-- node[lab,above,pos=0.5,yshift=1pt] {$p_2$} (cat);
  \draw[fl] (tr)  -- node[lab,above,pos=0.5,yshift=1pt] {$p_3$} (cat);

  \draw[fl] (cat) -- (ml);
  \draw[fl] (ml)  -- (out);

  % Bandeaux de niveau
  \node[lab,anchor=north] at (3.6,-3.05) {Niveau 0 --- modèles de base};
  \node[lab,anchor=north] at (10.9,-3.05) {Niveau 1 --- méta-apprenant};
  \draw[hypnogray!45,dashed,line width=0.5pt] (5.6,3.1) -- (5.6,-3.0);

\end{tikzpicture}
\vspace{0.4cm}
\caption[Principe de l'empilement de modèles]%
        {Principe de l'empilement : trois modèles hétérogènes produisent des vecteurs
         de probabilités, qu'un méta-apprenant apprend à combiner.}
\label{fig:stacking-principe}
\end{figure}

\subsection{Conditions d'efficacité et pièges}

L'empilement n'apporte un gain que sous deux conditions, et les ignorer conduit à des
résultats décevants ou, pire, illusoires.

La première est la \textbf{diversité des erreurs}. Si les modèles de base se trompent
sur les mêmes exemples, aucune combinaison ne peut les corriger. C'est l'argument qui
justifie de retenir des familles architecturales franchement distinctes plutôt que
plusieurs variantes d'un même réseau : un convolutif, un récurrent et un attentionnel
n'exploitent pas les mêmes régularités du signal, et n'échouent donc pas aux mêmes
endroits.

La seconde est l'\textbf{absence de fuite d'information vers le méta-apprenant}. Ce
point, identifié dès les premiers travaux sur le sujet~\cite{ting1999}, est le piège
classique de la méthode. Si le méta-apprenant est entraîné sur les prédictions que les
modèles de base produisent sur \emph{leurs propres données d'entraînement}, il apprend
sur des probabilités anormalement confiantes, puisque les modèles ont mémorisé ces
exemples. Il en déduit une pondération inadaptée aux données réelles. La parade
consiste à réserver une partition distincte, jamais vue par les modèles de base, pour
entraîner le méta-apprenant. C'est le protocole que nous suivrons, et le
chapitre~\ref{chap:realisation-ia} discutera dans quelle mesure il a effectivement été
respecté.

Un dernier point mérite attention. L'empilement doit être comparé non seulement au
meilleur modèle isolé, mais aussi au \textbf{vote moyen}, combinaison triviale qui ne
requiert aucun apprentissage. Un empilement qui ne bat pas la moyenne arithmétique des
probabilités n'apporte rien qu'une complexité supplémentaire. Cette comparaison, trop
souvent omise dans la littérature, figurera systématiquement dans nos résultats.

% =============================================================================
\section{Le dépistage du syndrome d'apnées obstructives du sommeil}
% =============================================================================

\subsection{Les instruments cliniques}

Le dépistage de première intention repose historiquement sur des questionnaires. Le
questionnaire de Berlin~\cite{netzer1999} et le score STOP-BANG~\cite{chung2008}
combinent des éléments d'anamnèse — ronflement, somnolence diurne, hypertension — et
des données anthropométriques. Leur profil est caractéristique : une sensibilité
élevée, qui les rend utiles pour ne pas manquer un patient à risque, mais une
spécificité faible, qui engendre de nombreux examens inutiles. Ils orientent, ils ne
concluent pas.

\subsection{Les approches fondées sur les signaux}

Une abondante littérature exploite les signaux enregistrés pendant la nuit. Une revue
systématique~\cite{mostafa2019} recense les approches par apprentissage profond
appliquées au flux nasal, aux sangles thoraciques, à l'électrocardiogramme et à
l'oxymétrie de pouls. Les performances rapportées sont solides.

Ces méthodes partagent cependant une contrainte : elles supposent disponibles les
capteurs respiratoires ou l'oxymétrie continue. Or ces canaux sont précisément ceux
qui alourdissent le montage, et ils sont absents des dispositifs d'enregistrement
allégés. Une méthode qui en dépend hérite de leur disponibilité.

\subsection{L'architecture du sommeil comme marqueur indirect}

Une troisième voie, moins explorée, consiste à exploiter la \emph{déformation de
l'architecture du sommeil} que produit le syndrome. Nous avons vu au
chapitre~\ref{chap:cadre} que les micro-éveils répétés fragmentent la nuit,
suppriment le sommeil profond et raccourcissent les épisodes paradoxaux. Cette
empreinte est lisible sur l'hypnogramme seul.

L'intérêt de cette voie est double. D'une part, l'hypnogramme est déjà produit par
l'étage de classification, ce qui rend la prédiction de sévérité gratuite en termes
d'acquisition. D'autre part, les marqueurs manipulés — efficacité du sommeil, durée
d'éveil intra-sommeil, index de fragmentation, latences — sont des grandeurs
cliniques familières au praticien, ce qui facilite considérablement l'explication de
la prédiction.

Le tableau~\ref{tab:etat-art-osa} situe ces familles d'approches.

\begin{table}[H]
\centering
\caption{Approches d'évaluation du risque de SAOS}
\label{tab:etat-art-osa}
\small
\begin{tabular}{@{}L{3.6cm}L{3.8cm}L{3.4cm}L{2.6cm}@{}}
\toprule
\textbf{Approche} & \textbf{Entrées} & \textbf{Atouts} & \textbf{Limites} \\
\midrule
Questionnaires\newline (Berlin, STOP-BANG) & Anamnèse, anthropométrie &
Immédiat, sans matériel & Spécificité faible, pas de sévérité \\
\addlinespace[3pt]
Oxymétrie nocturne & Saturation en oxygène & Matériel léger, bon dépistage &
Sévérité imprécise \\
\addlinespace[3pt]
Signaux respiratoires\newline et cardiaques & Flux nasal, sangles, ECG &
Performances élevées & Exige les capteurs respiratoires \\
\addlinespace[3pt]
\textbf{Architecture du sommeil}\newline \emph{(voie retenue)} & Hypnogramme, oxymétrie
sommaire, index d'éveil & Réutilise l'étage de staging ; marqueurs cliniquement
lisibles & Marqueur indirect \\
\bottomrule
\end{tabular}
\end{table}

% =============================================================================
\section{L'intelligence artificielle explicable}
% =============================================================================

\subsection{Une taxonomie des méthodes}

Le champ de l'intelligence artificielle explicable s'est structuré rapidement, au
point de justifier plusieurs synthèses~\cite{adadi2018,doshivelez2017}. Deux axes
suffisent à en organiser la lecture, et sont représentés à la
figure~\ref{fig:xai-taxo}.

Le premier axe oppose les méthodes \textbf{spécifiques} à une famille de modèles aux
méthodes \textbf{agnostiques}, applicables à n'importe quel prédicteur pourvu qu'on
puisse l'interroger. Le second oppose les explications \textbf{locales}, qui justifient
une prédiction individuelle, aux explications \textbf{globales}, qui décrivent le
comportement d'ensemble du modèle.

\begin{figure}[H]
\centering
\begin{tikzpicture}[
    x=1cm,y=1cm,
    racine/.style = {rectangle,rounded corners=2pt,draw=hypnoblue!65,fill=hypnolight,
                     align=center,font=\small\bfseries,minimum height=0.8cm,text width=3.6cm},
    branche/.style= {rectangle,rounded corners=2pt,draw=hypnoblue!50,fill=white,
                     align=center,font=\scriptsize\bfseries,minimum height=0.75cm,text width=3.0cm},
    feuille/.style= {rectangle,rounded corners=2pt,draw=hypnogray!45,fill=white,
                     anchor=west,align=left,font=\scriptsize,
                     minimum height=0.62cm,text width=3.35cm,inner xsep=5pt},
    ret/.style    = {rectangle,rounded corners=2pt,draw=hypnored!80,fill=hypnored!8,
                     anchor=west,align=left,font=\scriptsize\bfseries,
                     minimum height=0.62cm,text width=3.35cm,inner xsep=5pt},
    li/.style     = {hypnoblue!55,line width=0.7pt},
    bus/.style    = {hypnogray!55,line width=0.7pt},
  ]

  % ── Racine et deux branches ───────────────────────────────────────────────
  \node[racine]  (r)   at (7.4,3.55) {Méthodes d'explication};
  \node[branche] (spe) at (2.6,2.05) {Spécifiques\\à une architecture};
  \node[branche] (agn) at (10.4,2.05) {Agnostiques\\au modèle};

  \draw[li] (r.west)  -- ++(-1.2,0) |- (spe.north);
  \draw[li] (r.east)  -- ++(1.2,0)  |- (agn.north);

  % ── Feuilles : méthodes spécifiques ───────────────────────────────────────
  \node[feuille] (sal)  at (2.15, 0.95) {Cartes de saillance~\cite{simonyan2014}};
  \node[feuille] (gcam) at (2.15, 0.15) {Grad-CAM~\cite{selvaraju2017}};
  \node[feuille] (ig)   at (2.15,-0.65) {Gradients intégrés~\cite{sundararajan2017}};
  \node[feuille] (att)  at (2.15,-1.45) {Poids d'attention~\cite{bahdanau2015}};

  \draw[bus] (1.75,1.60) -- (1.75,-1.45);
  \foreach \n in {sal,gcam,ig,att}
    \draw[bus] (1.75,0 |- \n.west) -- (\n.west);
  \draw[bus] (1.75,1.60) -- (2.6,1.60) -- (2.6,1.68);

  % ── Feuilles : méthodes agnostiques ───────────────────────────────────────
  \node[feuille] (lime) at (9.95, 0.95) {LIME~\cite{ribeiro2016}};
  \node[ret]     (shap) at (9.95, 0.15) {SHAP~\cite{lundberg2017}};
  \node[feuille] (perm) at (9.95,-0.65) {Importance par permutation};
  \node[feuille] (pdp)  at (9.95,-1.45) {Dépendance partielle};

  \draw[bus] (9.55,1.60) -- (9.55,-1.45);
  \foreach \n in {lime,shap,perm,pdp}
    \draw[bus] (9.55,0 |- \n.west) -- (\n.west);
  \draw[bus] (9.55,1.60) -- (10.4,1.60) -- (10.4,1.68);

  \node[font=\scriptsize\itshape,text=hypnogray,anchor=north] at (7.0,-2.10)
       {la méthode encadrée en rouge est celle retenue dans ce travail};

\end{tikzpicture}
\vspace{0.3cm}
\caption[Taxonomie des méthodes d'explication]%
        {Taxonomie des principales méthodes d'explication, selon leur dépendance à
         l'architecture du modèle.}
\label{fig:xai-taxo}
\end{figure}

\subsection{Les méthodes spécifiques aux réseaux profonds}

Les \textbf{cartes de saillance}~\cite{simonyan2014} calculent le gradient de la sortie
par rapport à l'entrée : elles répondent à la question « quel échantillon du signal,
s'il changeait légèrement, modifierait le plus la décision ? ».
\textbf{Grad-CAM}~\cite{selvaraju2017} exploite les cartes d'activation des dernières
couches convolutives. Les \textbf{gradients intégrés}~\cite{sundararajan2017} corrigent
une faiblesse du gradient brut — la saturation — en intégrant le long d'un chemin
depuis une entrée de référence.

Les \textbf{poids d'attention} sont souvent présentés comme une explication native.
Cette lecture est contestée. Jain et Wallace~\cite{jain2019} montrent qu'il existe
fréquemment des distributions d'attention très différentes conduisant à la même
prédiction, ce qui interdit de tenir l'attention pour \emph{l'}explication. Wiegreffe
et Pinter~\cite{wiegreffe2019} nuancent en distinguant plusieurs acceptions du terme.
Le débat n'est pas tranché, et la prudence s'impose : les poids d'attention indiquent
où le modèle regarde, pas nécessairement pourquoi il décide.

\subsection{Les valeurs de Shapley et SHAP}

SHAP~\cite{lundberg2017} transpose à l'apprentissage automatique un concept issu de la
théorie des jeux coopératifs : la valeur de Shapley~\cite{shapley1953}. Le problème
originel consiste à répartir équitablement le gain d'une coalition entre ses membres.
Transposé, il devient : comment répartir l'écart entre la prédiction obtenue et la
prédiction moyenne, entre les variables d'entrée ?

La valeur de Shapley d'une variable $i$ est sa contribution marginale moyenne, calculée
sur tous les ordres d'arrivée possibles :

\begin{equation}
\phi_i = \sum_{S \subseteq F \setminus \{i\}}
         \frac{|S|!\,\bigl(|F| - |S| - 1\bigr)!}{|F|!}
         \Bigl[ f_{S \cup \{i\}}(x_{S \cup \{i\}}) - f_S(x_S) \Bigr]
\label{eq:shapley}
\end{equation}

\noindent où $F$ est l'ensemble des variables et $f_S$ le modèle restreint au
sous-ensemble $S$. Le terme entre crochets est l'apport de la variable $i$ lorsqu'elle
rejoint la coalition $S$ ; le coefficient pondère cet apport par la probabilité de
l'ordre d'arrivée correspondant.

Cette formulation confère à SHAP trois propriétés que peu de méthodes réunissent.
L'\textbf{additivité} garantit que la somme des contributions restitue exactement
l'écart à la valeur de base, ce qui autorise une lecture en cascade
(figure~\ref{fig:shap-cascade}). La \textbf{cohérence} assure qu'une variable dont
l'influence augmente dans un modèle modifié ne verra jamais sa contribution diminuer.
La \textbf{symétrie}, enfin, attribue la même valeur à deux variables strictement
interchangeables.

\begin{figure}[H]
\centering
\begin{tikzpicture}[
    x=1cm,y=1cm,
    barP/.style = {fill=hypnored!70,draw=hypnored,line width=0.4pt},
    barN/.style = {fill=hypnoblue!55,draw=hypnoblue,line width=0.4pt},
    et/.style   = {font=\scriptsize,anchor=east,text=hypnogray},
    va/.style   = {font=\scriptsize,text=hypnogray},
  ]

  % Axe de référence
  \draw[hypnogray!45,dashed,line width=0.6pt] (0,0.5) -- (0,-4.6);
  \node[font=\scriptsize\itshape,text=hypnogray,anchor=south] at (0,0.55)
       {valeur de base $E[f(x)]$};

  % Contributions successives
  \draw[barP] (0,-0.2)   rectangle (3.2,-0.65);
  \node[et] at (-0.15,-0.42) {index de fragmentation};
  \node[va,anchor=west] at (3.35,-0.42) {$+$};

  \draw[barP] (3.2,-1.0) rectangle (5.0,-1.45);
  \node[et] at (-0.15,-1.22) {\% du temps sous \SI{90}{\percent}};
  \node[va,anchor=west] at (5.15,-1.22) {$+$};

  \draw[barP] (5.0,-1.8) rectangle (6.1,-2.25);
  \node[et] at (-0.15,-2.02) {indice de masse corporelle};
  \node[va,anchor=west] at (6.25,-2.02) {$+$};

  \draw[barN] (5.3,-2.6) rectangle (6.1,-3.05);
  \node[et] at (-0.15,-2.82) {\% de sommeil profond};
  \node[va,anchor=west] at (6.25,-2.82) {$-$};

  \draw[barN] (4.9,-3.4) rectangle (5.3,-3.85);
  \node[et] at (-0.15,-3.62) {latence d'endormissement};
  \node[va,anchor=west] at (6.25,-3.62) {$-$};

  % Prédiction finale
  \draw[hypnored,line width=1pt,dashed] (4.9,-4.1) -- (4.9,-4.6);
  \node[font=\scriptsize\bfseries,text=hypnored,anchor=north] at (4.9,-4.65)
       {prédiction : sévère};

  \node[font=\scriptsize\itshape,text=hypnogray,anchor=west] at (7.0,-1.6)
       {\begin{minipage}{5.2cm}\raggedright
        Chaque barre est la contribution $\phi_i$ d'une variable.
        Les contributions rouges poussent vers une sévérité plus élevée,
        les bleues vers une sévérité moindre. Leur somme relie exactement
        la valeur de base à la prédiction.
        \end{minipage}};

\end{tikzpicture}
\vspace{0.2cm}
\caption[Lecture en cascade d'une explication SHAP]%
        {Lecture en cascade d'une explication SHAP : la propriété d'additivité relie
         la valeur de base à la prédiction individuelle.}
\label{fig:shap-cascade}
\end{figure}

Le calcul exact de l'équation~\eqref{eq:shapley} est exponentiel en le nombre de
variables et donc impraticable au-delà d'une vingtaine d'entrées. Lundberg
\emph{et al.}~\cite{lundberg2020} ont levé cet obstacle pour les modèles à base
d'arbres en proposant TreeSHAP, qui exploite la structure arborescente pour obtenir
les valeurs exactes en temps polynomial. C'est cette variante que nous emploierons,
ce qui explique un choix de conception détaillé au
chapitre~\ref{chap:realisation-ia} : l'explication est calculée sur un modèle à
arbres, même lorsque la prédiction est rendue par un ensemble hétérogène.

\subsection{Limites et objections}

Deux réserves doivent être formulées, car un mémoire qui présenterait l'explicabilité
comme une solution acquise manquerait de rigueur.

La première tient à la nature même de ces méthodes : elles expliquent le
\emph{modèle}, non le \emph{phénomène}. Une contribution SHAP élevée signale une forte
influence sur la prédiction, en aucun cas un lien de causalité physiologique. La
confusion est facile et lourde de conséquences en contexte médical.

La seconde est plus radicale. Rudin~\cite{rudin2019} soutient que, pour les décisions
à fort enjeu, il vaut mieux employer d'emblée un modèle intrinsèquement interprétable
que d'expliquer après coup une boîte noire — l'explication \emph{a posteriori} étant
par construction une approximation du comportement réel du modèle. L'argument est
sérieux et s'applique à notre travail. Nous y répondons par la conception : l'étage de
prédiction de sévérité repose sur des variables cliniques nommées et directement
interprétables, ce qui rapproche l'explication du raisonnement médical, et le modèle
sous-jacent est un ensemble d'arbres pour lequel TreeSHAP fournit des valeurs
\emph{exactes} et non approchées.

% =============================================================================
\section{Positionnement de notre contribution}
% =============================================================================

Le tableau~\ref{tab:positionnement} résume ce que cet état de l'art établit et ce
qu'il laisse ouvert.

\begin{table}[H]
\centering
\caption{Ce que l'état de l'art établit, ce que notre travail ajoute}
\label{tab:positionnement}
\small
\begin{tabular}{@{}L{3.2cm}L{5.3cm}L{5.3cm}@{}}
\toprule
\textbf{Domaine} & \textbf{Acquis de la littérature} & \textbf{Apport de ce travail} \\
\midrule
Classification\newline des stades &
Les architectures profondes atteignent l'accord inter-expert ; aucune famille ne
domine uniformément &
Évaluation systématique de trois familles sur deux montages et deux granularités,
sur une même cohorte et un même protocole \\
\addlinespace[3pt]
Apprentissage\newline ensembliste &
L'empilement de modèles hétérogènes réduit conjointement biais et variance &
Application au staging avec comparaison explicite au vote moyen, trop souvent omise \\
\addlinespace[3pt]
Prédiction\newline du SAOS &
Approches par questionnaire, oxymétrie ou signaux respiratoires &
Prédiction à quatre classes fondée sur l'architecture du sommeil, sans capteur
respiratoire \\
\addlinespace[3pt]
Explicabilité &
SHAP fournit des attributions additives et cohérentes &
Explication portée par des variables cliniques nommées, intégrée au compte rendu remis
au praticien \\
\addlinespace[3pt]
Mise en {\oe}uvre &
Les travaux publiés s'arrêtent au modèle &
Chaîne applicative complète : dossier patient, archivage, annotation, collaboration,
traçabilité \\
\bottomrule
\end{tabular}
\end{table}

% =============================================================================
\section{Conclusion}
% =============================================================================

Ce chapitre a parcouru les quatre corpus scientifiques mobilisés par le projet. La
classification automatique des stades est un problème mûr, où les performances
plafonnent au voisinage de l'accord inter-expert et où aucune famille architecturale
ne s'impose — constat qui justifie le recours à l'empilement de modèles hétérogènes.
L'apprentissage ensembliste offre le cadre théorique de cette combinaison, à condition
d'en respecter les deux exigences : diversité des erreurs et étanchéité entre les
données d'entraînement des deux niveaux.

Le dépistage du SAOS reste, lui, partagé entre des questionnaires peu spécifiques et
des méthodes exigeant des capteurs respiratoires. La voie de l'architecture du sommeil
apparaît sous-exploitée alors qu'elle se greffe naturellement sur un étage de
classification déjà présent.

Enfin, l'explicabilité dispose avec SHAP d'un socle théorique solide, dont nous avons
aussi rappelé les limites : ces méthodes expliquent le modèle et non le phénomène, et
l'objection de Rudin en faveur des modèles intrinsèquement interprétables reste
recevable. Notre réponse est architecturale plutôt que rhétorique, et se lira dans les
choix de conception.

Le chapitre suivant quitte l'état de l'art pour formaliser les besoins du système :
acteurs, exigences fonctionnelles et non fonctionnelles, \emph{Product Backlog} et
découpage du projet en sprints.
